\chapter{Opis projektu}

\section{Charakterystyka aplikacji}
Projektowana aplikacja to kompleksowy system informatyczny wspierający działalność nowoczesnej firmy cateringowej oferującej diety pudełkowe. Głównym celem systemu jest z jednej strony umożliwienie klientom łatwego i intuicyjnego zamawiania diet, a z drugiej – dostarczenie pracownikom firmy zaawansowanego panelu do zarządzania procesami biznesowymi. Aplikacja automatyzuje obsługę zamówień, zarządza magazynem i składnikami, a także wspiera planowanie produkcji posiłków oraz organizację logistyki i dostaw.

\section{Typy użytkowników}
W systemie wyróżniono dwie główne grupy użytkowników, dla których zaprojektowano odrębne interfejsy:
\begin{itemize}
    \item \textbf{Użytkownicy klienccy (Frontend)}: Dzielą się na użytkowników niezalogowanych (przeglądających ofertę i cenniki) oraz zalogowanych, którzy mogą składać zamówienia, modyfikować aktywne diety oraz zarządzać swoimi danymi i adresami dostaw.
    \item \textbf{Użytkownicy administracyjni (Backend)}: Obsługują panel zarządzania. Wyróżnia się tu administratora głównego (pełny dostęp), pracowników biura (obsługa zamówień), magazynierów (kontrola stanów składników) oraz logistyków/kurierów (obsługa tras).
\end{itemize}

\section{Funkcjonalności aplikacji}
Poniżej zestawiono kluczowe funkcjonalności systemu, podzielone na moduły klienckie i administracyjne.

\textbf{Funkcjonalności użytkownika (klienta):}
\begin{itemize}
    \item \textbf{Strona główna i przegląd diet}: Prezentacja oferty, sekcje promocyjne, filtrowanie diet (np. standardowa, wegetariańska, ketogeniczna) oraz szczegółowe opisy jadłospisów.
    \item \textbf{Konfigurator zamówienia}: Personalizacja diety (wybór kaloryczności, liczby dni, daty rozpoczęcia), dynamiczne przeliczanie ceny oraz wybór adresu dostawy.
    \item \textbf{Proces płatności}: Podsumowanie zamówienia i wybór metody płatności online.
    \item \textbf{Panel użytkownika i wsparcie}: Dostęp do historii zamówień, podgląd aktywnych diet, możliwość edycji danych osobowych oraz formularz kontaktowy i sekcja FAQ.
\end{itemize}

\textbf{Funkcjonalności administratora:}
\begin{itemize}
    \item \textbf{Dashboard (Panel główny)}: Widok podsumowujący liczbę zamówień, dostawy na dany dzień, generujący przychody oraz alerty magazynowe.
    \item \textbf{Zarządzanie zamówieniami i CRM}: Pełna obsługa zamówień (zmiana statusów: nowe, w realizacji, w dostawie), filtrowanie oraz baza klientów z historią zamówień i notatkami o alergiach.
    \item \textbf{Magazyn i Dostawcy}: Kontrola stanów magazynowych, historia zmian, alerty braków na podstawie systemu zapasu bezpieczeństwa oraz lista dostawców i cen zakupów.
    \item \textbf{Logistyka i Diety}: Planowanie tras dla kurierów, tworzenie diet, planowanie menu i przypisywanie składników z magazynu do konkretnych posiłków.
\end{itemize}